% =============================================================================
\chapter{Bilan}\label{chap:bilan}
% =============================================================================

\section{Principales difficultés}

Le projet a rencontré une quarantaine de difficultés, toutes documentées avec leur cause et leur
correction dans les chapitres précédents. Elles se regroupent en quatre familles :

\begin{table}[H]
\centering
\caption{Familles de difficultés et enseignements}\label{tab:difficultes}
\begin{tabularx}{\linewidth}{@{}P{3.3cm}P{5.2cm}L@{}}
\toprule
\textbf{Famille} & \textbf{Exemples} & \textbf{Enseignement} \\
\midrule
Environnement & ISO Ubuntu erronée, carte réseau e1000, absence d'iptables, VM éteintes après un redémarrage de l'hôte & Vérifier la plateforme réelle avant de dérouler la procédure prévue \\
Idempotence & PTR du DC réécrit à chaque passage, DNS \code{::1}, quotas réécrits, handler non exécuté & Un second passage est un test, pas une formalité : il révèle des conflits entre l'automatisation et le système \\
Vérifications trompeuses & Assertion réussie sur 0 cible, HTTP 200 avec un mauvais contenu, \code{--check} qui n'exécute pas les lectures, arrêt « TASK OK » d'une VM tuée & Une vérification doit échouer quand elle ne vérifie rien ; tester aussi le cas négatif \\
Mises à jour & Boucle Windows Update sur KB5007651, lecteur CD vu comme volume plein, collecteur apt imprécis & Croiser plusieurs sources (historique, journaux, état réel) avant de conclure \\
\bottomrule
\end{tabularx}
\end{table}

\section{Écarts au sujet}

\begin{itemize}
  \item \textbf{Proxmox VE 9.1 au lieu de VMware Workstation} — à faire valider par le formateur
        (avantages : snapshots, ajout de disque à chaud, installations et snapshots pilotables par API).
  \item Ressources légèrement supérieures au plan (vCPU) ; disque Windows en IDE (choix initial conservé).
  \item Ubuntu 24.04.\textbf{5}, dernière révision de la version LTS demandée.
\end{itemize}

\section{Fonctionnalités inachevées}

\begin{table}[H]
\centering
\caption{Points ouverts et action nécessaire}\label{tab:inacheves}
\begin{tabularx}{\linewidth}{@{}P{5cm}LP{2cm}@{}}
\toprule
\textbf{Point} & \textbf{Action nécessaire} & \textbf{Statut} \\
\midrule
Arrêt non propre de WIN-DC01 (INC-2026-10-23-01) & Validation de l'équipe, application de \code{ShutdownWithoutLogon = 1}, test d'arrêt réel & \partiel \\
Bilans individuels (séance 4) et revue croisée & À rédiger par chaque membre & \afaire \\
Contributions de Kenzi Bendjelloul dans l'historique Git & Commits à ajouter par l'intéressé & \afaire \\
Notification des alertes & Ajouter Alertmanager (messagerie de l'équipe) & \afaire \\
KB5007651 & Suivi manuel, retrait de l'exclusion quand Microsoft aura corrigé & \afaire \\
Démarrage automatique des VM & Décider de l'option \code{onboot} sur l'hôte & \afaire \\
\bottomrule
\end{tabularx}
\end{table}

\section{Perspectives}

\begin{itemize}
  \item \textbf{Supervision hors de l'hôte} : une sonde externe (ou un second hôte) pour détecter un arrêt
        complet du laboratoire, angle mort mis en évidence par l'incident UC-05.
  \item \textbf{Chiffrement des flux d'administration} : WinRM en HTTPS, TLS pour Prometheus et Grafana.
  \item \textbf{Centralisation des journaux} (Loki et Grafana Alloy) pour corréler plus vite les événements
        de plusieurs serveurs.
  \item \textbf{Alerte prédictive sur le disque}, à valider après une semaine de fonctionnement continu.
  \item \textbf{Intégration d'Ubuntu au domaine Active Directory} pour centraliser les identités.
\end{itemize}

\section{Bilan de l'équipe}

Le projet a atteint ses objectifs : un laboratoire hybride entièrement décrit par un dépôt Git, des
configurations rejouables et idempotentes, des secrets chiffrés, une supervision qui détecte à la
fois l'indisponibilité d'un serveur et celle d'un service, et une chaîne de livraison où chaque
changement est validé automatiquement puis relu. La règle « une preuve par résultat, y compris les
échecs » a structuré le travail : la plupart des défauts corrigés ont été trouvés en vérifiant
qu'une vérification pouvait échouer.

\info{Bilans individuels}{%
Les bilans individuels demandés par les supports des séances 4 et 6 (opération maîtrisée, commande
la plus utile, difficulté isolée, risque identifié, reproductibilité) sont à rédiger par chaque
membre de l'équipe ; ils ne sont pas rédigés à leur place dans ce rapport.}

\section*{Conclusion}
\addcontentsline{toc}{section}{Conclusion}

InfraReady disposait au départ de serveurs administrés à la main. Elle dispose désormais d'une
infrastructure de laboratoire dont l'état est reproductible depuis le dépôt, dont la maintenance
s'arrête d'elle-même si les conditions ne sont pas réunies, et dont les anomalies sont visibles avant
les utilisateurs. Les écarts et limites sont identifiés et accompagnés d'une action ; le principal
point ouvert, l'arrêt non propre du contrôleur de domaine, est compris, reproduit et sa correction
est prête. Ce socle peut être étendu sans remise en cause, dans l'ordre des perspectives ci-dessus.
